\section{可解群}

设$G$为单群。因为换位子群$G^{(1)}$是$G$的正规子群，单群的定义给出
\[
G^{(1)}=\{e\}\quad\text{或}\quad G^{(1)}=G.
\]
若$G^{(1)}=\{e\}$，则$G$中任意两个元素$a,b$的换位子都等于$e$，即$a^{-1}b^{-1}ab=e$，从而$ab=ba$。所以$G$是交换群。这说明，非交换单群必满足
\[
G^{(1)}=G.
\]

对于一般的群$G$，换位子群$G^{(1)}$也是正规子群，但它既可能等于$\{e\}$，也可能等于$G$，还可能是非平凡真子群。继续对$G^{(1)}$取换位子群，便得到下面的递降子群列。

\begin{definition}
令$G^{(0)}=G$，递归定义
\[
G^{(k)}=\bigl(G^{(k-1)}\bigr)^{(1)}\qquad(k\geq1).
\]
所得子群列
\[
G=G^{(0)}\geq G^{(1)}\geq G^{(2)}\geq\cdots
\]
称为$G$的\textbf{导出列}。其中每一项都是前一项的换位子群，因而是前一项的正规子群。
\end{definition}

若$G$为有限群，导出列不能无限次严格下降，因而只有两种可能：一是到达$\{e\}$，此后各项都等于$\{e\}$；二是在某个非平凡子群处稳定，即存在$k$使
\[
G^{(k)}=G^{(k+1)}=\cdots\ne\{e\}.
\]
这里，一旦相邻两项相等，继续取换位子群便得到后面的各项也相等。

\begin{definition}
若存在正整数$k$使$G^{(k)}=\{e\}$，则称$G$为\textbf{可解群}。
\end{definition}

交换群都是可解群，因为交换群的换位子群等于$\{e\}$。非交换单群的导出列则一直等于自身，所以非交换单群不可解。

\begin{example}
对称群$S_n$在$n\leq4$时可解，在$n\geq5$时不可解。分别计算如下。

\textbf{$S_1,S_2$：}$S_1$为平凡群，$S_2=\{e,(12)\}$为二阶循环群，都是交换群，因而可解。

\textbf{$S_3$：}由$S_3/A_3\cong C_2$为交换群，得$S_3^{(1)}\leq A_3$。另一方面，按照$[a,b]=a^{-1}b^{-1}ab$的约定，
\[
[(12),(23)]=(12)(23)(12)(23)=(132).
\]
因此$S_3^{(1)}$包含$(132)$生成的整个$A_3$，故
\[
S_3^{(1)}=A_3=\{e,(123),(132)\}.
\]
由于$A_3$是三阶循环群，它是交换群，所以$S_3^{(2)}=A_3^{(1)}=\{e\}$。

\textbf{$S_4$：}同理由$S_4/A_4$交换得$S_4^{(1)}\leq A_4$，且上述换位子说明$S_4^{(1)}$含有一个3-轮换。$S_4^{(1)}$在$S_4$中正规，而所有3-轮换在$S_4$中互相共轭，所以$S_4^{(1)}$包含全部3-轮换。由3-轮换生成$A_4$，得
\[
S_4^{(1)}=A_4.
\]
令
\[
V_4=\{e,(12)(34),(13)(24),(14)(23)\}.
\]
由$A_4/V_4\cong C_3$为交换群，得$A_4^{(1)}\leq V_4$。又直接计算有
\[
[(123),(124)]=(132)(142)(123)(124)=(12)(34).
\]
把$(12)(34)$分别用$(123)$及$(132)$共轭，得到
\begin{align*}
(123)(12)(34)(132)&=(14)(23),\\
(132)(12)(34)(123)&=(13)(24).
\end{align*}
因$A_4^{(1)}\trianglelefteq A_4$，它包含上述三个双对换，故$A_4^{(1)}=V_4$。$V_4$是交换群，所以
\[
S_4^{(1)}=A_4,\qquad S_4^{(2)}=V_4,\qquad S_4^{(3)}=\{e\}.
\]

\textbf{$n\geq5$：}任取$\sigma,\tau\in S_n$，其换位子为偶置换，因为
\[
\operatorname{sgn}(\sigma^{-1}\tau^{-1}\sigma\tau)
=\operatorname{sgn}(\sigma)^{-1}\operatorname{sgn}(\tau)^{-1}
\operatorname{sgn}(\sigma)\operatorname{sgn}(\tau)=1.
\]
所以$S_n^{(1)}\leq A_n$。同时$S_n^{(1)}\trianglelefteq S_n$，故它也是$A_n$的正规子群。上述$[(12),(23)]=(132)$表明$S_n^{(1)}\ne\{e\}$。由$A_n$为单群，得到
\[
S_n^{(1)}=A_n.
\]
又因$A_n$是非交换单群，$A_n^{(1)}=A_n$，于是
\[
S_n^{(k)}=A_n\ne\{e\}\qquad(k\geq1).
\]
因此$S_n$不可解，$A_n$也不可解。
\end{example}

若$G$可解，则其导出列经过有限步到达$\{e\}$。每个$G^{(i)}$都是$G^{(i-1)}$的换位子群，故相应的商群
\[
G^{(i-1)}/G^{(i)}
\]
均为交换群。因此可解群具有一条相邻两项正规、相应商群交换的有限子群列。下面证明，这一性质也足以保证可解。

\begin{theorem}\label{p2:solvable-series}
群$G$可解，当且仅当存在有限递降子群列
\[
G=G_0\geq G_1\geq\cdots\geq G_s=\{e\},
\]
其中$G_i\trianglelefteq G_{i-1}$，且$G_{i-1}/G_i$为交换群，$i=1,\ldots,s$。
\end{theorem}
\begin{proof}
\textbf{必要性：}若$G^{(s)}=\{e\}$，取$G_i=G^{(i)}$。由导出列的定义，每一项在前一项中正规，且每个商群$G^{(i-1)}/G^{(i)}$交换，故满足条件。

\textbf{充分性：}设存在上述子群列。由定理\ref{p2:abelianization}，对任一正规子群$N\trianglelefteq G$，有
\[
G/N\text{为交换群}\quad\Longleftrightarrow\quad G^{(1)}\leq N.
\]
因此，$G_0/G_1$交换首先给出$G^{(1)}\leq G_1$。下面对$k$归纳，证明
\[
G^{(k)}\leq G_k\qquad(1\leq k\leq s).
\]
若已知$G^{(k-1)}\leq G_{k-1}$，则$G^{(k-1)}$中元素的换位子也是$G_{k-1}$中元素的换位子，所以
\[
G^{(k)}=\bigl(G^{(k-1)}\bigr)^{(1)}\leq G_{k-1}^{(1)}.
\]
另一方面，$G_{k-1}/G_k$交换给出$G_{k-1}^{(1)}\leq G_k$。合并即得$G^{(k)}\leq G_k$，完成归纳。

最后取$k=s$，有$G^{(s)}\leq G_s=\{e\}$，故$G^{(s)}=\{e\}$，即$G$可解。
\end{proof}

进一步讨论有限可解群时，需要说明何时还能在相邻两项之间插入新的正规子群。

\begin{definition}
设$N$为$G$的真正规子群。若不存在$G$的正规子群$N_1$使
\[
N<N_1<G,
\]
则称$N$为$G$的\textbf{极大正规子群}。这里要求中间子群$N_1$在$G$中正规。
\end{definition}

\begin{proposition}
设$N\trianglelefteq G$且$N\ne G$。商群$G/N$为单群，当且仅当$N$为$G$的极大正规子群。
\end{proposition}
\begin{proof}
由正规子群的对应定理，$G/N$中的正规子群与$G$中包含$N$的正规子群一一对应，其中$\{N\}$对应$N$，$G/N$对应$G$。

因此，$G/N$有非平凡真正规子群，当且仅当存在$N_1\trianglelefteq G$满足$N<N_1<G$。否定两边，便得到所述等价关系。
\end{proof}

\begin{theorem}\label{p2:finite-solvable-series}
有限群$G$可解，当且仅当存在有限严格递降子群列
\[
G=H_0>H_1>\cdots>H_t=\{e\},
\]
其中$H_i\trianglelefteq H_{i-1}$，且每个商群$H_{i-1}/H_i$都是素数阶循环群。对平凡群，允许$t=0$。
\end{theorem}
\begin{proof}
\textbf{必要性：}设$G$为有限可解群。由定理\ref{p2:solvable-series}，可取一条相邻商群均交换的有限子群列，并删去重复项，写成
\[
G=G_0>G_1>\cdots>G_k=\{e\}.
\]
若某个商群$G_{i-1}/G_i$不是单群，则它有非平凡真正规子群$L$。由对应定理，$L$的完全反像$G_i'$满足
\[
G_i<G_i'<G_{i-1},\qquad G_i'\trianglelefteq G_{i-1}.
\]
又因$G_i\trianglelefteq G_{i-1}$，也有$G_i\trianglelefteq G_i'$。因此可在原群列中插入$G_i'$。

插入后新增的两个商群满足
\[
G_i'/G_i\cong L,\qquad
G_{i-1}/G_i'\cong(G_{i-1}/G_i)/(G_i'/G_i).
\]
前一个是交换群$G_{i-1}/G_i$的子群，后一个同构于该交换群的商群，因此两者都仍为交换群。

只要还有一个相邻商群不是单群，就重复上述插入过程。由于$G$有限，严格递降子群列的长度有界，所以不能无限次插入。最终得到一条不能再插入的子群列
\[
G=H_0>H_1>\cdots>H_t=\{e\},
\]
其中每个商群$H_{i-1}/H_i$都既交换又单。

交换单群必为素数阶循环群：任取非单位元$a$，其生成的非平凡子群$\langle a\rangle$因交换而正规，所以等于整个群；若这个循环群的阶为合数，则又有非平凡真子群，与单性矛盾。因此上述各商群均为素数阶循环群。

\textbf{充分性：}若存在所述子群列，则每个相邻商群都为循环群，因而交换。由定理\ref{p2:solvable-series}，$G$可解。
\end{proof}

\begin{remark}
有限可解群的上述子群列不一定唯一。若
\[
|H_{i-1}/H_i|=p_i\qquad(i=1,\ldots,t),
\]
其中$p_i$均为素数，则逐次使用$|H_{i-1}|=p_i|H_i|$，得到
\[
|G|=p_1p_2\cdots p_t.
\]
由整数的素因子分解唯一性，$t$等于$|G|$的素因子个数，重复的素因子重复计算。因此虽然子群列未必唯一，它的长度以及相邻商群的素数阶（允许调整顺序）都是由$|G|$唯一确定的。
\end{remark}

\begin{proposition}\label{p2:solvable-extension}
可解群的子群、同态像均可解。若$N\trianglelefteq G$，且$N$与$G/N$均可解，则$G$可解。
\end{proposition}
\begin{proof}[补充证明]
\textbf{子群：}设$H\leq G$。$H$中元素的换位子也是$G$中元素的换位子，故$H^{(1)}\leq G^{(1)}$。逐次取换位子群并归纳，得到
\[
H^{(i)}\leq G^{(i)}\qquad(i\geq0).
\]
若$G^{(k)}=\{e\}$，则$H^{(k)}=\{e\}$，所以$H$可解。

\textbf{同态像：}设$\varphi:G\to Q$为满同态。由
\[
\varphi([a,b])=[\varphi(a),\varphi(b)]
\]
可知$\varphi(G^{(1)})\leq Q^{(1)}$。反过来，因$\varphi$满射，$Q$中任一换位子都可如此得到；这些换位子生成$Q^{(1)}$，故$\varphi(G^{(1)})=Q^{(1)}$。反复应用于各级导出子群，得
\[
\varphi(G^{(i)})=Q^{(i)}.
\]
所以$G^{(k)}=\{e\}$蕴含$Q^{(k)}=\{e\}$，即$Q$可解。

\textbf{由正规子群与商群推出原群可解：}设$N^{(s)}=\{e\}$、$(G/N)^{(r)}=\{N\}$。对自然满同态$\pi:G\to G/N$，上面的等式给出
\[
\pi(G^{(r)})=(G/N)^{(r)}=\{N\},
\]
因而$G^{(r)}\leq\ker\pi=N$。再逐次取$s$次换位子群，得到
\[
G^{(r+s)}=\bigl(G^{(r)}\bigr)^{(s)}\leq N^{(s)}=\{e\}.
\]
因此$G$可解。
\end{proof}

